\chapter{Introduction}
In recent years, there has been a growing concern about the increasing number of fatalities caused by car accidents resulting from driver distraction. According to the WHO, driver distraction is a major contributing factor to road traffic accidents worldwide. Globally, it's estimated that 20-30\% of all road traffic accidents are caused by driver distraction (World Health Organization, 2011), In the European Union, it's estimated that driver distraction is a contributing factor in 10-30\% of all road accidents (European Transport Safety Council, 2020). Therefore distracted driving has become a leading cause of accidents, accounting for a significant portion of road fatalities worldwide. This alarming trend has necessitated the development of advanced algorithms capable of recognizing driver actions and behaviors to prevent such accidents and promote road safety.

To address this critical issue, there has been a surge in the development of action recognition systems specifically tailored for the car interior. These systems aim to identify and analyze various driver actions, such as hand gestures, head movements, and body postures, in real-time. By accurately recognizing these actions, potential distractions and unsafe behaviors can be identified, allowing for timely intervention and accident prevention.

Traditional supervised learning methods have been widely employed for action recognition tasks. However, they heavily rely on annotated datasets, which are costly and time-consuming to produce. Annotated datasets require manual labeling of each action instance, posing practical limitations when it comes to scaling up the training process.

To overcome these challenges, self-supervised learning has emerged as a promising alternative approach. Unlike traditional supervised learning, self-supervised learning leverages unannotated data to learn useful representations. By formulating the learning task as a pretext task, such as predicting missing parts or ordering sequences, self-supervised learning allows the model to learn meaningful representations without the need for explicit annotations. This approach has gained significant attention due to its ability to utilize large amounts of readily available unannotated data.

Furthermore, the increasing development of autonomous vehicles has amplified the demand for accurate and robust human action recognition systems. As autonomous cars become more prevalent, it becomes crucial to understand and interpret human actions within the vehicle environment. Accurate recognition of driver actions and behaviors enables autonomous systems to respond appropriately, ensuring passenger safety and efficient vehicle control.

In light of these challenges and opportunities, this master thesis aims to enhance vehicle interior action recognition by leveraging contrastive self-supervised learning techniques with 3D human skeleton representations. By utilizing unannotated data, this approach can effectively learn informative representations that capture the nuances of driver actions in the car interior. The proposed method will contribute to advancing the field of action recognition by providing a scalable and efficient solution to enhance driver safety in both conventional and autonomous vehicles.

Through extensive experimentation and evaluation, this thesis seeks to demonstrate the effectiveness of contrastive self-supervised learning with 3D human skeleton representations in improving the accuracy and robustness of vehicle interior action recognition systems. The findings of this research will not only contribute to the development of safer driving environments but also provide valuable insights into the application of self-supervised learning techniques in the context of using sequence data as human skeleton data.

Overall, this master thesis addresses the need for accurate and efficient action recognition systems in the car interior to prevent accidents caused by driver distraction. The adoption of contrastive self-supervised learning with 3D human skeleton representations represents a novel and promising approach to achieving this objective. By combining the advancements in self-supervised learning and human action recognition, this research aims to enhance driver safety and realize reliable autonomous vehicle systems.

\section{Problem Statement}
Human action recognition has been extensively studied using different modalities, such as RGB videos or depth cameras. However, these modalities often come with significant computational costs, making them less suitable for resource-constrained environments. In response to this challenge, human skeleton data has emerged as a promising alternative. Skeleton data offers a more efficient representation of human actions, reducing the computational complexity and storage requirements, and also they provide a lightweight and efficient alternative that not only reduces computational overhead but also offers the potential to generalize across domains, enabling broader applicability in real-world scenarios, as explained in \cite{martin2023viewpoint}. By focusing on the underlying structure and motion of the human body, skeleton data provides a compact representation that is computationally less expensive while still capturing the essential information for action recognition.

Traditional supervised techniques for action recognition heavily rely on annotated datasets. However, the process of manually annotating data for action recognition is time-consuming and labor-intensive, limiting its scalability for large-scale applications. In this context, self-supervised learning has gained traction as a promising solution. It enables models to learn meaningful representations without the need for explicit annotations. This approach addresses the challenge of costly and time-consuming data annotation, making it a more efficient and scalable solution for action recognition.

While self-supervised learning has shown promise in various domains, its application to skeleton data for action recognition is still in its early stages. Further research and investigation are necessary to explore the full potential of self-supervised learning in this context. The thesis aims to bridge this gap by investigating the effectiveness of self-supervised learning techniques for action recognition using human skeleton data. By leveraging the inherent structure and motion information in the skeleton data, the thesis seeks to enhance the accuracy and robustness of action recognition models, even with limited or no labeled data.

Furthermore, the lack of action recognition datasets specifically dedicated to driver behavior inside the car presents another challenge. Existing datasets often focus on general action recognition but fail to capture the nuanced actions and behaviors specific to drivers within the vehicle environment. To address this gap, the thesis will focus on the test and evaluation of an action recognition dataset tailored specifically to driver behavior inside the car. This dedicated dataset will incorporate a wide range of driver actions and behaviors, providing a realistic and comprehensive evaluation platform for action recognition algorithms in the automotive context.

In summary, this master thesis aims to tackle the challenges associated with human action recognition by leveraging the benefits of skeleton data and self-supervised learning. The research will investigate the feasibility and effectiveness of self-supervised learning techniques for action recognition using skeleton data. Additionally, the thesis will contribute to domain adaptation from traditional human action recognition to driver action recognition inside the car, enabling accurate evaluation and benchmarking of action recognition algorithms in the automotive domain.

\section{Research Objectives}
The first research objective is to design and develop a contrastive self-supervised learning approach specifically tailored for action recognition on 3D human skeleton data. The focus will be on leveraging the benefits of contrastive learning to learn rich and discriminative representations from unannotated skeleton data. The proposed approach will explore different contrastive learning strategies in which it relies on relevant data augmentations that can optimize the action recognition performance. The effectiveness of the developed approach will be rigorously evaluated and benchmarked against state-of-the-art methods in the field of action recognition on skeleton data.

The second research objective is to evaluate the proposed contrastive self-supervised learning approach on datasets dedicated to capturing driver behavior inside the car, encompassing a comprehensive range of driver actions and behaviors, in which these datasets were created and developed by Fraunhofer IOSB - where the master thesis is conducted -and KIT.  The evaluation will focus on assessing the performance of the proposed method in accurately recognizing driver actions within the unique context of the vehicle environment. The evaluation process will involve extensive quantitative analysis and comparison against existing approaches, providing insights into the effectiveness of the proposed method for driver behavior recognition in practical applications.

The third research objective aims to delve deeper into the landscape of contrastive self-supervised learning techniques. This involves investigating the strengths and limitations of various techniques, analyzing their effects on different aspects of action recognition, and identifying potential areas of improvement.

By achieving these research objectives, this master thesis aims to contribute to the advancement of action recognition techniques on 3D human skeleton data and the evaluation on dedicated driver behavior datasets will provide valuable insights for real-world deployment in automotive applications. The results obtained from this research will help pave the way for improved driver safety and facilitate the development of intelligent systems in autonomous driving and driver assistance technologies.

In summary, the research objectives of the master thesis can be defined as follows:
\begin{itemize}
	\item Design and implementation of a contrastive self-supervised learning approach for action recognition using 3D human skeleton representations.
	\item Training and evaluating of the approach on publicly available datasets for 3D skeleton data-based action recognition (e.g. NTU RGB+D, Drive\&Act).
        \item  Analyze the impact of different contrastive self-supervised learning techniques on various evaluation protocols of action recognition.
.
\end{itemize}


\section{Thesis Outline}
The master thesis will follow a structured approach to address the research objectives outlined in the previous sections. The thesis will be organized into the following chapters:

\textbf{Chapter 1:} Literature review
will provide a comprehensive review of the relevant literature in the field of action recognition using different modalities. The review will delve into action recognition with RGB data, and action recognition with a skeleton data highlighting the various approaches employed, including RNNs, CNNs, GCNs, and transformers based. The chapter will also focus on the concept of self-supervised learning (SSL) and its relevance in action recognition tasks. Specifically, it will discuss the pretext tasks and contrastive learning methods commonly used in self-supervised learning approaches. The literature review will offer a theoretical foundation and critical analysis of existing methods and highlight their strengths, weaknesses, and limitations.

\textbf{Chapter 2:} The methodology chapter will detail the proposed approach for self-supervised action recognition on skeleton data. It will cover the preparation and processing of skeleton data, including any required normalization or data augmentation techniques. The chapter will present an overview of different approaches and architectures considered for the task, incorporating insights from the literature review. The training strategy, including the choice of a specific SSL and contrastive learning frameworks, will be described.

\textbf{Chapter 3:} Experiments chapter will showcase the experimental design and evaluation protocols for assessing the performance of the model will also be presented. The datasets used in the experiments, including any data preprocessing specific to each dataset, will be discussed.

\textbf{Chapter 4:} Results and Evaluation
will present the results of the experiments conducted to evaluate the proposed contrastive self-supervised learning approach. The evaluation metrics used to measure the performance of the model will be discussed, highlighting the chosen metrics and their relevance to action recognition tasks. Ablation studies will be conducted to analyze the contribution of different components or variations of the proposed approach. The chapter will also include a comprehensive comparison with state-of-the-art methods, showcasing the strengths and advantages of the proposed approach.

\textbf{Chapter 5:}
The final chapter will summarize the findings and contributions of the master thesis. It will provide a concise overview of the research objectives, methodology, and the key results obtained. The chapter will discuss the implications and significance of the findings in the context of action recognition on 3D human skeleton data. Additionally, any limitations or potential areas for future research will be highlighted. The conclusion chapter will serve as a reflection on the achievements of the thesis and provide closure to the overall research journey.